% META: {"id":"1NSI-TC-EX-035","chapitre":"1NSI-TYPES-CONSTRUITS","type_objet":"exercice","capacites":["1NSI-TYPES-CONSTRUITS-C4"],"parcours":1,"competences":["analyser"],"duree_min":5,"mode_creation":"ex_nihilo","sources_inspiration":[],"status":"approved","origin":"needs_review","status_history":["needs_review","audited","accepted","approved"],"release_acceptance":"RELEASE_OWNER_BATCH_ACCEPTANCE","acceptance_closure_digest":"sha256:9b3ccf9a81c5520fbb7e03b7b2d3e7bf2057a02834b6d908bf3be5f49f3b553f","acceptance_receipt_digest":"sha256:4fad86f0282e0fa4e0419cca1ad6379ae7af5a280c04a4a90880f90a1c044242"}
% BEGIN-TRACE
% eleve = {"nom": "Alice", "classe": "1NSI", "moyenne": 14.5}
% print(eleve["nom"])
% print(eleve["moyenne"])
% print(len(eleve))
% EXPECTED
% Alice
% 14.5
% 3
% END-TRACE
\begin{exercice}{1NSI-TC-EX-035}{1}{5}
On donne le dictionnaire suivant :
\begin{python}
eleve = {"nom": "Alice", "classe": "1NSI", "moyenne": 14.5}
print(eleve["nom"])
print(eleve["moyenne"])
print(len(eleve))
\end{python}
\begin{enumerate}
  \item Donner, sans exécuter, les trois affichages.
  \item Quelles sont les clés de ce dictionnaire ? Quelles sont les valeurs ?
  \item Écrire l'instruction qui ajoute la clé \lstinline{"email"} avec la valeur
        \lstinline{"alice@lycee.fr"}.
\end{enumerate}
\end{exercice}
